\section{The structure theorem}

The goal of this section is to prove \Cref{thm:maintechnical} (see Section~\ref{sec:introduction} for the statement). We start by giving the algorithm for the $c$-bounded case, that is the proof of Theorem~\ref{thm:subexpalgo}, and refocusing our efforts on acyclic instances.

\subsection{A subexponential algorithm}
\label{sec:dag}

In order to prove \Cref{thm:subexpalgo}, we will utilize the following result.

\begin{theorem}[Lemma 2.1 of~\cite{ChitnisFHM20}]\label{thm:smalldeletion_scss}
Suppose that $H$ is an edge-minimal solution for the \textsc{Strongly Connected Steiner Network} problem on the instance $(G,T)$, where $|T|=k$. Then there is a set $W\subseteq V(H)$ of $9k$ vertices such that deleting $W$ from $H$ results in a graph of treewidth at most $2$.
\end{theorem}

We start by generalizing the above theorem to our setting.

\begin{lemma}\label{lem:smalldeletion}
Let $H$ be an edge-minimal solution to \Dsnfull where $\calD$ is a $c$-bounded class. Then there is a set of $(30c+29)k$ vertices in $V(H)$ whose deletion from $H$ results in a graph $H'$ of treewidth at most $2$.
\end{lemma}

\begin{proof}
Let $V_1,\dots,V_t$ be the vertex sets of the strongly connected components of $H$ that are not singletons. The degree of $V_i$, denoted by $d(V_i)$, is the number of edges in $H$ with one endpoint in $V_i$ and one endpoint outside $V_i$.

We claim that if $V_i$ has no terminals, then  $d(V_i)\geq 3$. Suppose the contrary: $V_i\cap T=\emptyset$ and $d(V_i)\leq 2$. Notice that $d(V_i)\neq 1$ as in such a case we could remove all edges of $H$ induced by $V_i$ as well as the edge that enters/exits $V_i$. Suppose now that $d(V_i)=2$. Notice that the two edges that have one endpoint in $V_i$ cannot both be entering or both be exiting $V_i$ as again that would be redundant in $H$. Thus $V_i$ has one edge entering at some vertex $u\in V_i$ and one edge exiting from some vertex $v\in V_i$. Let $\pi$ be a path in $H[V_i]$ from $u$ to $v$. Notice that since $|V_i|\geq 2$ and $H[V_i]$ is strongly connected, there must be some edge $e$ in $H[V_i]$ that is not on the path~$\pi$. However all terminal-to-terminal paths passing through $V_i$ can be realized using $\pi$, since if such a path enters $V_i$ then it enters at $u$ and exits at $v$. Consequently, the edge $e$ is redundant, contradicting the minimality of $H$.

We say that a vertex $v$ is a \emph{portal} of $V_i$ if it is an endpoint of some edge $e\in E(H)$ that enters or exits $V_i$. Let $P$ denote the set of all portals, and let $P_i$ be the set of portals in $V_i$. Consider now a strongly connected subgraph $H[V_i]$ that contains at least one terminal, and let $T_i=V_i\cap T$. Since $H[V_i]$ is strongly connected, it is an edge-minimal solution of \scssfull for the terminal set $T_i\cup P_i$: indeed, if some edge $e$ can be removed from $H[V_i]$ to maintain strong connectivity on $T_i\cup P_i$, then $H-e$ would also be feasible for the original problem, which contradicts the edge-minimality of $H$. We can therefore apply \Cref{thm:smalldeletion_scss} for the graph $H[V_i]$ and terminal set $T_i\cup P_i$: there is a set $X_i$ of at most $9|T_i\cup P_i|$ vertices such that $H[V_i\setminus X_i]$ has treewidth at most $2$.

We now delete the following vertices from $H$:
\begin{itemize}
\item the portal set $P$
\item all the sets $X_i$
\item the set $Y$ of vertices outside $\bigcup_i V_i$ that have degree at least $3$.
\end{itemize}
Let $H^*$ be the resulting graph. Each connected component of $H^*$ is either a subgraph of some strongly connected component $H[V_i]$, or it lies outside $\bigcup_i V_i$, thus its vertices have maximum degree at most $2$: in both cases, it has treewidth at most $2$. Consequently, $\tw(H^*)\leq 2$.

In order to bound the number of deleted vertices, let us first introduce an acyclic instance of \Dsnfull based on $(G,T,D)$ and $H$: this is required so that we can use the $c$-boundedness of $\calD$ . Let $G/(\bigcup V_i)$ be the graph obtained by contracting each vertex set $V_i$ into a single vertex (deleting all loops and parallel edges). Notice that the picture $H'$ of $H$ under this contraction is acyclic. In the demand graph $D$, we identify each vertex set $T_i$ individually, resulting in a demand graph $D'$ and terminal set $T'$. Notice that since $\calD$ is closed under identification, we have $D'\in \calD$. 

We claim that $H'$ is an edge-minimal solution to the 
\Dsnfull instance $(G'=H',T',D')$ (where $H'$ and $D'$ are acyclic). First we show that $H'$ is feasible. If $D'$ has an edge $u'v'$, then the edge has an ancestor $uv$ in $D$, so $H$ contains some path $P_{uv}$ connecting the terminals $u$ and $v$ (where $u$ and $v$ cannot be located in the same set $V_i$). When we apply the contraction on $P_{uv}$, we get a path $P'$ connecting $u'$ and $v'$ in $H'$, concluding the proof of feasibility for $H'$.

To prove the edge-minimality of $H'$, suppose the contrary: that there is some edge $a'b'$ in $H'$ whose deletion does not break feasibility. Let $ab$ be an edge of $H$ that contracts to $a'b'$. Notice that $ab$ is not induced by any component $V_i$. We claim that $H-ab$ is feasible for $(G,T,D)$. Consider a demand edge $uv\in E(D)$. If the terminals $u$ and $v$ are in the same $V_i$, then they are connected within $H[V_i]$ and thus they are also connected in $H-ab$. Otherwise these terminals are identified with some distinct terminals $u'$ and $v'$, respectively, and $D'$ has a demand $u'v'$. Thus $H'-a'b'$ connects $u'$ and $v'$ via some path $P'$. Let $u$ and $v$ be arbitrary vertices of $G$ that contract to (or are equal to) $u'$ and $v'$, respectively.

Using $P'$, we will now build a path $P$ connecting $u$ and $v$ in $H-ab$. First, we pick for each edge $x'y' \in P'$ an arbitrary edge whose image after contraction is $x'y'$. Consider now an edge pair $x'y', y'z'$ on $P'$: if $y'$ comes from the contraction of $V_i$, then we may have to select $xy_1$ as the ancestor of $x'y'$ and $y_2z$ as the ancestor of $y'z'$, where $y_1 \neq y_2$ are different portals. But since $(H-ab)[V_i]$ is strongly connected, there is a path $P_i$ connecting the portals $y_1$ to $y_2$ in $(H-ab)[V_i]$. We concatenate this path with the edges $xy_1$ and $y_2z$. Using the same technique, we can build a path from some $u_2$ to some $v_2$ where either $u_2=u$ or $u_2$ and $u$ are in the same component $V_j$ (and the analogous statement holds for $v$ and $v_2$). We can again use the strong connectivity of the components $V_i$ to get a path $P$ from $u$ to $v$. Consequently, $H-ab$ is feasible for the original problem, which contradicts the edge-minimality of $H$.

Finally, we note that for any pair $V_i,V_j$ with $i\neq j$ there can be at most one edge going between $V_i$ and $V_j$. Since these are strongly connected components, there cannot be two edges in different directions, as that would make them into a single strongly connected component. Suppose that $uv$ and $u'v'$ both go from $V_i$ to $V_j$. Then either of these edges can be removed without affecting feasibility, which contradicts the edge-minimality of $H$. This property implies that in the proposed contraction the degree of the vertex that we get from contracting $V_i$ is equal to $d(V_i)$.

We can now bound the size of the deleted set. Since $\calD$ is $c$-bounded, we have that $\sum_{v\in V(H')} d^*(v)\leq ck$. Since each non-0 term in the sum comes from some contraction of $V_i$ or a deleted vertex from $Y$, we obtain the following.
\begin{equation}\label{eq:cbounded}
\sum_{v\in V(H')} d^*(v) = \sum_i d^*(V_i) + \sum_{v\in Y} d^*(v)\leq ck,\end{equation}
where $d^*(V_i)=\max(0,d(V_i)-2)$.
Recall that if $V_i$ has degree at most $2$, then it must have at least one terminal. It follows that it can have at most $2$ portals, and thus $|P_i|\leq 2|T_i|$. If it has degree at least three, then it has at most $d(V_i)= d^*(V_i)+2\leq 3d^*(V_i)$ portals.

Thus the total number of deleted vertices can be bounded as:
\begin{align*}
|Y|+|P|+\Big|\bigcup_i X_i\setminus P\Big|&\leq |Y|+|P|+\sum_i 9|T_i\cup P_i|\\
&\leq |Y|+9|T|+10|P|\\
&= |Y|+9k + 10\cdot \left(\sum_{i\,:\,d(V_i)\geq 3} |P_i| + \sum_{i\,:\,d(V_i)\leq 2} |P_i|\right)\\
&\leq |Y| + 9k + 10\cdot \left(\sum_{i} 3d^*(V_i) + \sum_{i\,:\,d(V_i)\leq 2} 2|T_i|\right)\\
&\leq (30c+29)k,
\end{align*}
where the second inequality uses that the sets $V_i$ are disjoint, and the last inequality uses the bound~\eqref{eq:cbounded}.
Consequently, we have removed at most $(30c+29)k$ vertices from $\bigcup_i V_i$, which concludes the proof.
\end{proof}

The next lemma can essentially be found within the proof of Lemma 2.2 in~\cite{ChitnisFHM20}. 
We reproduce the proof for completeness.

\begin{lemma}\label{lem:twboundfromdeletion}
If $G$ is a planar graph where deleting $k\geq 1$ vertices results in a graph of treewidth $w\geq 1$, then $\tw(G)\leq 3w\sqrt{k}$.
\end{lemma}

\begin{proof}
By the planar grid theorem~\cite{RobertsonST94}, there is a constant $\bar c$ such that any planar graph of treewidth $\bar c \omega$ has a grid minor of size $\omega \times \omega$. If $G$ has treewidth at least $\bar c \cdot \lceil 3w\sqrt{k}\rceil$, then it has a grid minor of size at least $\lceil 3w\sqrt{k}\rceil$. This minor can be decomposed into $\left\lfloor \frac{ \lceil 3w\sqrt{k}\rceil}{w+1}\right\rfloor \cdot \left\lfloor \frac{ \lceil 3w\sqrt{k}\rceil}{w+1}\right\rfloor$ vertex disjoint grid minors, each of size at least $(w+1)\times(w+1)$. Notice that $\left\lfloor \frac{ \lceil 3w\sqrt{k}\rceil}{w+1}\right\rfloor \cdot \left\lfloor \frac{ \lceil 3w\sqrt{k}\rceil}{w+1}\right\rfloor\geq k+1$ for $k,w\geq 1$ integers, thus there is at least one $(w+1)\times(w+1)$ grid minor where no vertex has been deleted. This intact grid minor has treewidth at least $w+1$, which contradicts our assumptions.
\end{proof}

We are now ready to prove \Cref{thm:subexpalgo}.

\begin{proof}[Proof of \Cref{thm:subexpalgo}]
Consider an input $(G,T,D)$. By \Cref{lem:smalldeletion}, we can find a set of $(30c+29)k$ vertices in $H$ whose deletion results in a graph of treewidth at most $2$. \Cref{lem:twboundfromdeletion} then bounds the treewidth of $H$ as $\tw(H)\leq 3\cdot 2 \cdot \sqrt{(30c+29)k}< 47\sqrt{ck}=O(\sqrt{k})$.
This concludes the proof.
\end{proof}